% =====================================================================
% PLANTILLA: Documento general (taller, guía, notas de clase, ...)
% Área de Matemáticas -- María Mercedes Cantillo · Colegio San Alberto Magno
% ---------------------------------------------------------------------
% 1. Copia este archivo con otro nombre (p. ej. taller-integrales.tex) en
%    esta misma carpeta (el estilo y el logo se toman de ../estilo/).
% 2. Cambia los datos de la sección "DATOS DEL DOCUMENTO".
% 3. Usa \encabezado  si el estudiante debe escribir su nombre (taller),
%    o  \encabezado* si es material de lectura (notas de clase).
% 4. Compila con:  latexmk -pdf taller-integrales.tex
% =====================================================================
\documentclass[11pt]{article}
% El estilo y el logo viven en la carpeta compartida plantillas/estilo/
\makeatletter\def\input@path{{../estilo/}}\makeatother
\usepackage{mmcantillo}

% ========================= DATOS DEL DOCUMENTO ========================
\tipodocumento{Taller}        % Taller, Guía, Notas de clase, ...
\asignatura{Cálculo}          % Cálculo, Física, Trigonometría, ...
\titulo{La integral definida como área}
\grado{11°}
\periodo{IV}
% \anio{2026}                 % por defecto: el año actual
% ======================================================================

\begin{document}

\encabezado

\seccion{Conceptos clave}

\begin{definicion}[Integral definida]
Sea $f$ continua en $[a,b]$. La integral definida de $f$ entre $a$ y $b$ es
\[
  \int_a^b f(x)\dd x \;=\; \lim_{n\to\infty} \sum_{i=1}^{n} f(x_i^{*})\,\Delta x,
  \qquad \Delta x = \frac{b-a}{n}.
\]
\end{definicion}

\begin{teorema}[Teorema fundamental del cálculo]
Si $f$ es continua en $[a,b]$ y $F$ es una antiderivada de $f$, entonces
\[
  \int_a^b f(x)\dd x = F(b) - F(a).
\]
\end{teorema}

\begin{ejemplo}[Área entre dos curvas]
El área entre $y = x$ y $y = x^2$ en $[0,1]$ es
\[
  A = \int_0^1 \left(x - x^2\right)\dd x
    = \eval{\left(\frac{x^2}{2} - \frac{x^3}{3}\right)}{0}{1}
    = \frac12 - \frac13 = \frac16 .
\]

\begin{center}
\begin{tikzpicture}
\begin{axis}[mmc, xmin=-0.2, xmax=1.4, ymin=-0.2, ymax=1.4,
             xtick={0,0.5,1}, ytick={0.5,1},
             width=0.5\linewidth, height=0.42\linewidth,
             legend pos=outer north east, legend style={draw=none, font=\footnotesize}]
  \addplot+[name path=f, domain=0:1.3] {x};
  \addplot+[name path=g, domain=0:1.2] {x^2};
  \addplot[fill=black!15, draw=none] fill between[of=f and g, soft clip={domain=0:1}];
  \legend{$y=x$, $y=x^2$}
\end{axis}
\end{tikzpicture}
\end{center}
\end{ejemplo}

\begin{nota}
Si $f(x) < 0$ en parte del intervalo, la integral cuenta esa región con
signo negativo. Para hallar el \emph{área} se integra $\abs{f(x)}$.
\end{nota}

\seccion{Ejercicios}

\begin{preguntas}

\pregunta Calcula $\displaystyle\int_1^3 \left(2x + 1\right)\dd x$ y
verifica el resultado con el área de un trapecio.
\espacio{3cm}

\pregunta Halla el área de la región limitada por $y = 4 - x^2$ y el eje $x$.
\espacio{3cm}

\pregunta Explica con tus palabras por qué $\displaystyle\int_{-\pi}^{\pi}
\sen x \dd x = 0$.
\renglones{3}

\end{preguntas}

\end{document}
